\section{Introduction}
\label{sec:intro}

Multimodal language models have become genuinely good at visual work. They caption and answer
questions about natural images, read documents, charts and tables, transcribe text in photographs,
follow video, and drive graphical interfaces from screenshots alone. Over the last two years the
direction of progress on these tasks has been steeply upward, and on several of them the reported
numbers sit near careful human annotators. TODO: Rewrite this.

It does not follow that these models reason about spatial reasong (TODO: Fisx this with geom
reasoning etc. ). Most of the competence arrives through one
architecture: a vision encoder trained to align images with text, and an adapter that projects its
output into the language model's token stream. What that encoder is rewarded for preserving is
semantic. It answers \emph{what is in this picture}, and it answers it well, because that is what
the training objective asks of it. Metric and relational structure is not rewarded and need not
survive the projection: exact angles, exact incidences, which of two nearly identical shapes lies
in front, how many times one line crosses another. Competence at recognising a scene and
competence at reasoning about its geometry are different, and an adapter trained for the first
gives no guarantee of the second. When the task moves to geometric and spatial reasoning, the
ranking of models by their image-task scores need not survive. (fix the wordings).

% ============================================================================
% STAR EXAMPLE, NOT YET VERIFIED. DO NOT SUBMIT WITH THIS PARAGRAPH AS IS.
% The intended example is the "Vision Language Models Are Blind" line of work,
% which reports that frontier models scoring well on hard visual question
% answering fail at elementary geometric perception: whether two circles
% overlap, counting intersections of two line plots, counting nested squares.
% Recorded reference: Rahmanzadehgervi, Bolton, Taesiri and Nguyen, ACCV 2024,
% arXiv:2407.06581.
% VERIFY authors, venue and the specific task list against the arXiv record
% before this enters the paper. If any part does not hold exactly as stated,
% CUT the example rather than soften it. See editorial rule 7 in Appendix C.
% ============================================================================

Why Origami for Spatial and Geometric reasoning

(DELETE AND WRITE AGAIN)
Origami is an unusually clean place to look for the missing capability. It is an art form whose
entire content is geometric: a folder works from lines, reflections and incidences on a single
uncut sheet, and every decision is constrained by every decision before it. (Mention about the
spatial thing - using the fact about the layer ordering etc. -  ) It also demands
spatial reasoning of a specific kind, because the sheet stops being flat the moment folding
begins. The folder has to track where each piece of paper now sits, which pieces lie above which,
and which of them are still joined to each other through the original sheet. A crease pattern, the
flat record left behind when a finished model is unfolded, contains all of that history and none
of its order. Reading a sequence back out of it is the skill that separates a practitioner from
someone who can recognise a crane.

We use origami as an instrument. \textbf{The model is given two things: the crease pattern, and the
final folded state. It is asked for the sequence of folds that turns the flat sheet into that
state.} Both endpoints are supplied and the path between them is not, which is the whole of the
task. The pattern says which lines were creased at some point during the folding but says nothing
about the order they were creased in, and the final state says where the paper ended up but not how
it got there.

 -  need a graph here to illustrate

This is trivial to generate and hard to solve, which is what makes it usable for evaluation.
Folding forward is a single pass of an engine that records what it did. Deciding whether a crease
pattern is reachable by simple folds at all is NP hard \citep{arkin-map,akitaya-hard}, so recovering
the sequence that produced one is at least as hard. The same engine that generates a sample
therefore rules on any proposed step of a solution exactly, using an operation it already
implements, with no learned component in the loop, no tolerance to tune and no oracle to trust.

That asymmetry is what a benchmark wants, and it is what spatial and physical reasoning has
generally lacked. Benchmarks built on static images or multiple choice never execute the model's
proposal, so an answer that is right for the wrong reason is indistinguishable from one that is
right. Benchmarks built on learned surrogate simulators do execute the proposal, but against an
approximation whose wrong verdicts look exactly like its right ones, which leaves the benchmark's
own error rate unknown and unmeasurable from inside. Human judgement is exact in a sense and does
not scale. The consequence is worth stating plainly, because it motivates everything below:
\textbf{a benchmark can only be as trustworthy as the judge inside it.} Here the judge is the
generator run forwards, so difficulty is controllable rather than sampled, ground truth is recorded
rather than annotated, and negative examples are unlimited, because an illegal fold is produced by
asking for one.

Two recent benchmarks make origami a target for multimodal evaluation, and the difference between
them and this work is the difference between asking a model a question and executing its answer.
\textbf{GamiBench} \citep{gamibench} pairs crease patterns with folded shapes rendered from six
viewpoints and scores multiple-choice visual question answering, including a set of deliberately
impossible patterns a model ought to reject. \textbf{OrigamiSpace} \citep{origamispace} instruments
each instance with a pattern diagram, the folding process and a final image, and defines four tasks
over them, from pattern prediction to crease-pattern code generation. In both, a proposal is graded
by comparison against a stored target: the model selects among rendered alternatives, or emits an
artifact that is matched to a reference. Neither steps a paper-exact engine and asks whether
\emph{this} move, on \emph{this} stack, is something paper could do. Both therefore measure
recognition of folded geometry, which is a real capability, rather than the ability to search for a
sequence that produces it. Section~\ref{sec:related-gap} makes the comparison in full.

The distinction matters because the search is where the difficulty lives. Determining the layer
ordering of a flat folding is NP hard on its own, even when a valid mountain-valley assignment is
supplied \citep{bern-hayes}. No amount of pattern recognition substitutes for search on a problem
of that shape, and a benchmark that never executes a proposed step cannot tell a model that
searched from a model that recognised, nor report how much search either one needed.

This is also why the benchmark reports a deterministic search baseline alongside the models rather
than only against them. The stack deepens with every fold, the enumerated action set grows with it,
and the depth at which exhaustive search stops being affordable is a property of the corpus that can
be measured rather than asserted. \texttt{[TO RUN: the deterministic breadth-first baseline of
Section~\ref{sec:experiments}, reported per stratum, giving the depth at which it stops solving and
the states it expands. Until that number exists, no claim is made here about where blind search
fails.]} What the benchmark is designed to reward is the use of geometry: looking at the crease
pattern and the final state and inferring which fold could plausibly have been last, then working
backwards. \textbf{Geometry is supplied to the model as input and is not the thing being scored.
Being able to act on it is.} A model that cannot read a reflection off a pattern has no way to
choose a direction other than enumeration, and the baseline is what says how far enumeration alone
gets.

The task is also hard in ways a maze or a grid-world is not. It gets harder as it proceeds: every
fold thickens the stack and constrains what the next fold may do, so the difficulty of step $k$
depends on steps $1$ through $k-1$. And a solver has to maintain two kinds of state that constrain
each other. The crease graph is fixed once the pattern is given, but vertex coordinates are
reflected at every step, the overlap structure changes at every step, and the layer ordering
changes at every step and is \emph{chosen} rather than computed. Determining that ordering is NP
hard even when a valid mountain-valley assignment is supplied \citep{bern-hayes}. Mazes have
geometry alone; here an independent combinatorial layer is tracked alongside it. The action space
is not invented for the occasion either: simple folds have been defined and their complexity
studied for two decades \citep{arkin-map}, so the operations an agent is given have names and
claims about them can be checked against a literature.

We contribute the following.

\begin{enumerate}
\item \textbf{A dataset, and the generator that built it.} CP2Seq holds 600 origami samples, each
  pairing a crease pattern with the fold sequence that produced it, the folded state after every
  step, and per-sample difficulty metrics. Ground truth is constructed rather than annotated,
  because the sequence is recorded as the fold happens. Samples are generated under two named
  action models from the simple folding literature and stratified by fold depth and by coupling.
  The corpus ships as a generator and a manifest rather than as data: every sample rebuilds
  byte-identically from its seed, so the dataset can be regenerated, extended, or re-stratified at
  a different difficulty by anyone who wants a different corpus than ours.

\item \textbf{An evaluation framework and environment, with progressive tiers of assistance.} The
  model is given the crease pattern and the final state and must recover the sequence, so it works
  from partial information about the structure throughout. The environment it works against is also
  the verifier: stepping it is the same act as asking for a verdict. It executes one candidate fold
  and either returns the resulting state or refuses with a named reason, such as the sheet would
  tear or the fold creases nothing. It performs no search and makes no choices, so proposing,
  pruning and backtracking stay with the model. Around it sits a graded tool belt, from a full belt
  with rendered views down through withheld vision and raw pass/fail to no tools at all, so the
  contribution of each assistance level is measured rather than assumed. The framework also
  supplies the scoring protocol: answers are judged by replaying them through the engine, at two
  declared levels of equality -- (edited distance and CP) - -, because one crease pattern admits
  many valid folded states and the
  recorded sequence is not guaranteed to be the shortest.

\item \textbf{An evaluation of frontier multimodal models} against this dataset through this
  framework, run across the progressive tiers of tool assistance and with a memorization control on
  anonymized geometry. \texttt{[TO RUN]}

\item \textbf{A measured scope boundary.} 89.3\%% fact:probeC.provenNotPct
  of the 366% fact:corpus.instagram.total
  real crease patterns in Flat-Folder's \texttt{examples/instagram/} corpus \citep{flatfolder}
  provably lie outside all-layers simple folding. The benchmark's action space is bounded by
  evidence rather than by assertion, and the boundary is reported rather than buried.
\end{enumerate}

We find that the task is beyond current models at the setting we could afford to run. ((No model
solved a single sample past the easy stratum  - ---are you sure, then we should give a graph, shows
what's newest model performance))), and on the easy stratum a deterministic breadth-first
search over exactly the action set the models are given solves 54.7\% against 20.4\% for the
best-covered model. Blind search beats every model tested, on the only stratum where anything
succeeds at all. The dominant failure is not exhausting the search budget but revisiting states
already seen: 48\% of that model's attempts terminate in detected state cycling.

\textbf{Caveat.} Every model attempt
ran at low reasoning effort, so the supported claim is about models at that setting rather than
about the frontier in general; Section~\ref{sec:results-threats} states this and
Section~\ref{sec:limitations} carries it as a limitation.- -----put in appendix

